% =============================================================================
% WORKING NOTES (not for submission) --- internal memo
% Zonality of SOC turnover: storyline development + diagnostic investigation
% Session 2026-06-25 (Lorenzo + Claude). Standalone-compilable.
% =============================================================================
\documentclass[11pt,a4paper]{article}
\usepackage[utf8]{inputenc}
\usepackage[margin=2.4cm]{geometry}
\usepackage{amsmath}
\usepackage{booktabs}
\usepackage{xcolor}
\usepackage{hyperref}
\hypersetup{colorlinks=true,linkcolor=black,urlcolor=blue!60!black}
\usepackage{enumitem}
\setlength{\parskip}{0.4em}

\title{\textbf{Working notes: the zonality of SOC turnover}\\[2pt]
\large Storyline decision and diagnostic investigation of the prediction maps}
\author{Internal memo --- L.\ Menichetti (with Claude Code)}
\date{2026-06-25 \quad\textcolor{red}{\small DRAFT / NOT FOR CIRCULATION}}

\begin{document}
\maketitle

\begin{center}\fbox{\parbox{0.92\textwidth}{\small\textbf{\textcolor{red}{SUPERSEDED
2026-06-26 --- historical record only.}} Going-forward planning has moved to
\texttt{storyline\_outline\_2026-06-26.\{tex,pdf\}} (the story bible) and
\texttt{results\_reconciliation\_2026-06-26.md} (number map + figure dispositions
+ manuscript checklist). This file is the 2026-06-25/26 exploration log; it is no
longer maintained. Do not edit further.}}\end{center}

\medskip
\noindent\textbf{Purpose.} Capture the state of the discussion so we can resume
cleanly. This documents (i)~the chosen storyline, (ii)~what we learned trying to
build a global ``zonality'' map, (iii)~a diagnostic investigation of the
prediction pipeline (including two false alarms and one real issue), and
(iv)~where the story now stands and the next concrete step.

% =============================================================================
\section{Storyline decision}
% =============================================================================

We commit to \textbf{framing~A}: the \emph{zonality of SOC decay} as an
\emph{interpretive framework}, presented as the main contribution. The argument:
SOC persistence is an emergent ecosystem property (Schmidt et al.\ 2011); if so, transit times
($\tau$) carry a spatially structured residual \emph{beyond climate}, analogous
to classical pedogenic zonality. The deliverable that turns this from slogan into
tool is an observation-constrained \emph{map} of that beyond-climate modulation.

Supporting roles: the global $\tau$ map and CMIP6/ESM contrast are the
empirical payload (B); the edaphic/mycorrhizal mechanism is the interpretive
texture (C). The ESM comparison is \emph{repositioned} from ``validation'' to
``motivating illustration'' (ESMs concentrate $\tau$ variation in a polar
extremum and lack mid-latitude structure); this removes the dependence of the
headline on the depth/pool comparability caveat, since the load-bearing claim
becomes how spatial variation is \emph{organised}, not its absolute magnitude.

% =============================================================================
\section{Building the zonality map: what works and what fails}
% =============================================================================

We must define ``zone'' \emph{from the data}. We tested several constructions and
scored each by how much of the \textbf{observed climate-residual} of
$\ln\tau$ they explain (point level, out-of-bag; $\eta^2$ = between-zone variance
fraction). Residual variance $\approx 0.32$ (on the biological subset).

\begin{table}[h]
\centering
\begin{tabular}{lcc}
\toprule
Construction & explains observed climate-residual & verdict\\
\midrule
Continuous RF modulation (M7$-$M1) & \textbf{13.4\%} (ceiling) & validated signal\\
\quad discretised, 5 zones & 10.5\% & \textbf{works}\\
\quad discretised, 8 zones & 11.5\% & works\\
argmax ``dominant process'' & $\sim$0\% & fails (collinear channels)\\
single hard-rule tree (CART) & $\sim$0\% & fails (signal diffuse)\\
$k$-means on raw predictors & 0\% & fails\\
\bottomrule
\end{tabular}
\caption{Only a \emph{supervised} construction works: discretising the validated
continuous RF ``non-climate modulation field'' recovers most of the achievable
signal at 5--8 zones. A Köppen-style ``which process dominates here'' map is
\emph{not} supported.}
\end{table}

\noindent\textbf{Why the intuitive map fails.} The three non-climate channels
(Edaphic, Land~Use, Biological prediction shifts) are spatially collinear:
a PCA of the signed $[E,L,B]$ adjustment vector gives PC1~$=67\%$ with near-equal
positive loadings ($0.59,0.53,0.60$; pairwise $r=0.43$--$0.61$). The
beyond-climate signal is therefore largely a \emph{single shared modulation axis}
(``how strongly non-climate processes bend $\tau$ away from the climate
expectation''), not three independent process-geographies. This is a cleaner,
more defensible framing than a process-mosaic --- and it matches the data.

\noindent\textbf{Recommended deliverable (3 layers):}
\begin{enumerate}[nosep]
  \item continuous global \emph{modulation field} = primary zonality-intensity map;
  \item a discrete 5--8 zone companion (quantile bins of the field) for communication;
  \item per-zone driver characterisation (predictor profiles / ALE) restoring the
        ``which process where'' reading as an honest \emph{secondary} layer.
\end{enumerate}

% =============================================================================
\section{Diagnostic investigation of the prediction maps}
% =============================================================================

While preparing to build the map we stress-tested the global prediction pipeline
(script~14). Net conclusion: \textbf{the maps are sound; the only real problem is
coverage.} Two alarms were chased down and \emph{refuted}.

\begin{table}[h]
\centering
\begin{tabular}{lc}
\toprule
Hypothesis & verdict\\
\midrule
M1 climate raster is corrupted (mean 35.6 vs point 21.8\,yr) & \textcolor{red}{refuted} --- coverage artifact\\
SoilGrids / covariate scaling bug in script~14 & \textcolor{red}{refuted} --- all 19 predictors $r\!\approx\!1.00$\\
Raster $(M7-M1)$ does not track the signal ($-0.15$) & \textcolor{red}{refuted} --- in-sample extraction artifact\\
\textbf{Microbial layers cover only 29\% of land} & \textcolor{blue}{\textbf{confirmed --- the real issue}}\\
\bottomrule
\end{tabular}
\caption{Outcome of the pipeline stress-test.}
\end{table}

\begin{itemize}[nosep]
  \item \textbf{No M1 scaling bug.} The 35.6 vs 22.3\,yr gap between M1 and M7
  is a \emph{domain} effect: M1 restricted to M7's footprint $=23.8 \approx
  M7=22.3$. All six climate input rasters match the training ranges; offsets in
  central tendency are the temperate \emph{sampling bias} of the training points.
  \item \textbf{No covariate scaling bug.} Per-predictor training-vs-raster check
  at 5{,}000 points: all 19 predictors (climate, SoilGrids edaphic with
  \texttt{d\_factor}, terrain, microbial) have $r\approx1.00$, median ratio
  $\approx1.00$.
  \item \textbf{The $-0.15$ ``inconsistency'' was our own error}: comparing
  observations to \emph{in-sample} raster predictions at the training points the
  models were fit on yields a meaningless in-bag residual. The correct
  out-of-bag comparison gives the non-climate adjustment correlating
  \textbf{$+0.37$} with the observed climate-residual. \emph{Any earlier
  raster-based $\eta^2\!\approx\!0$ numbers were contaminated by this same
  in-sample baseline; the point-level OOB results above stand.}
\end{itemize}

\noindent\textbf{The real issue --- coverage.} M1 (climate) covers
2{,}152{,}154 land pixels; the microbial-dependent models (M4/M6/M7) cover only
\textbf{631{,}252 = 29\% of land} --- the microbial layers (SPUN, Barceló,
fungal proportion) are \texttt{NA} over $\sim$71\% of land (high latitudes,
drylands). So the ``global'' $\tau$ map, the ESM comparison, and the planned
zonality field inherit these holes precisely where the ESM-divergence story is
strongest (boreal/Arctic). This is a data gap, not a code bug.

% =============================================================================
\section{Option 2: M5 as the global-coverage workhorse}
% =============================================================================

To get genuine global coverage we adopt \textbf{M5 (Climate $+$ Edaphic $+$
Land~Use)} --- no microbial dependency --- as the main map, with M7 retained as a
biologically-enriched overlay where microbial data exist. The M5 modulation field
$\ln(M5/M1)$ is saved at
\texttt{outputs/MRT\_predictions/zonality\_modulation\_M5\_M1\_logratio.tif}
(mean $0.018$, SD $0.279$, range $[-1.47, 2.01]$).

\begin{table}[h]
\centering
\begin{tabular}{lccc}
\toprule
Model & land coverage & $\Delta R^2$ vs climate (OOB) & explains climate-residual\\
\midrule
M5 (climate+edaphic+landuse) & \textbf{49\%} (all soil-bearing land) & 0.017 & 3.7\%\\
M7 (+ biological)            & 29\% & 0.063 & 13.4\%\\
\bottomrule
\end{tabular}
\caption{Coverage improves (29\%$\to$49\%) but the abiotic signal is weak.}
\end{table}

\noindent\textbf{Key emergent finding.} The \emph{biological} (mycorrhizal) group
carries most of the beyond-climate zonality, yet it is the least globally mapped.
\emph{The strongest beyond-climate signal is the least mappable.} Reframed
positively, this is a result in its own right --- soil biology is the dominant
non-climate control on SOC turnover, and the absence of global soil-biology maps
is the limiting factor. It motivates a concrete call for global microbial mapping
and supports a \textbf{dual-layer deliverable}: a global \emph{abiotic} zonality
field (M5) plus a \emph{biological enrichment} overlay (M7) where data permit.

% =============================================================================
\section{Resolving process classes without a per-pixel ``which-where'' map}
% =============================================================================

The channel collinearity blocks \emph{one} thing, not everything. Keep two
questions apart:
\begin{itemize}[nosep]
  \item \textbf{Q1 --- ``which process class dominates \emph{where}?''}
  (a per-pixel attribution map). \textcolor{red}{Blocked.} Channels correlate
  $r\approx0.4$--$0.6$, enough to make a pixel-by-pixel argmax unstable (the
  salt-and-pepper map). A Köppen-style process-territory map is \emph{not}
  supportable.
  \item \textbf{Q2 --- ``which process \emph{class} carries beyond-climate
  signal, and how much?''} (aggregate). \textcolor{blue}{Works.}
\end{itemize}

\noindent\textbf{Key reconciliation: correlated $\neq$ redundant.}
Correlation $\sim0.5$ is enough to break per-pixel argmax (Q1 fails) but
\emph{not} enough to make classes redundant: biology still adds large
\emph{unique} variance (M7 reaches 13.4\% of the residual vs M5's 3.7\%).
``Collinear'' was too strong --- the channels are correlated, not redundant.

\noindent\textbf{Resolve Q2 by nested increments, not 3-way attribution.}
Use the nested ladder $M1\!\to\!M5\!\to\!M7$. Each increment is a \emph{single}
comparison, so its magnitude \emph{and} geography are interpretable (unlike a
simultaneous 3-way argmax):
\begin{itemize}[nosep]
  \item $M5-M1$ = \textbf{abiotic} modulation (edaphic + land use; land use is
  redundant, so $\approx$ edaphic). Globally mappable ($\sim$49\%/all soil).
  \item $M7-M5$ = \textbf{biological enrichment} (the \emph{unique} biological
  contribution given abiotic + climate). Substantial and largely independent
  (since $M7\gg M5$), on the 29\% footprint.
\end{itemize}
Caveat: nested $\Delta R^2$ is order-dependent (``biology given abiotic''); a
commonality analysis (shared vs unique) is the rigorous version, but the
headline (biology adds much beyond abiotic; land use redundant) is robust.

\noindent\textbf{Coverage two levels = Tier~2.} The coverage limit is not a
bolt-on caveat; it \emph{is} the scope boundary between the nested tiers:
\textbf{Tier~1} (a beyond-climate zonality exists and is spatially organised)
is the overarching thesis and does \emph{not} depend on the 29\% --- it is
visible in M5 over all soil-bearing land. \textbf{Tier~2} resolves that
zonality into a global abiotic level ($M5-M1$) and a biological-enrichment
zoom-in ($M7-M5$, 29\%). Narrative: a global Level-1 abiotic zonality, then a
Level-2 deep-dive where soil biology is mapped. ``Biological conclusions hold
within the microbially-mapped domain'' + a call for global soil-biology maps.

% =============================================================================
\section{Routes to a clustered zone map}
% =============================================================================

Two distinct jobs: \textbf{(i) discover} the zones (which pixels group together)
and \textbf{(ii) describe} them with rules (portable, interpretable). Three
constructions, increasing in richness. \emph{Terminology: call the classes
\textbf{zones / turnover regimes}, not ``syndromes''.}

\begin{itemize}[nosep]
  \item \textbf{Route A --- cluster the nested-increment space}
  $[M1,\,M5\!-\!M1,\,M7\!-\!M5]$. Cheap internal sanity check only; uses
  existing rasters. Collinearity-safe (increments, not raw channels).
  \item \textbf{Route C (alone) --- regression tree on the modulation field.}
  One step: leaves are zones, splits are rules. Zones defined by \emph{level} of
  modulation (intensity bands). Transparent, portable, Köppen-like.
  \item \textbf{Route B --- cluster per-pixel driver contributions (TreeSHAP).}
  Each pixel gets a $\pm$-years budget per predictor; clustering groups pixels
  shaped by the same driver combination. Richest for \emph{insight}, but
  attribution has its own collinearity caveat and is not portable (needs the RF
  + SHAP to classify a new site). An analysis, not a tool.
\end{itemize}

\noindent\textbf{The B$\to$C bridge.} B and C are not a forced chain --- B
\emph{discovers}, C \emph{describes}. To connect them, take B's per-pixel
cluster labels and fit a classification tree whose \emph{response is the cluster
ID} and whose inputs are the predictors. That tree emits a rule per zone (``if
sand $>$ 60\% $\to$ zone~3; else if fungal $>$ 0.2 \& pH $<$ 5 $\to$
zone~5''). Why bother over C-alone: a regression tree merges regions with the
same $\tau$-shift even if driven differently (a dryland shortened by sand vs by
fungal\,+\,low-pH both read ``$-5$\,yr'' $\to$ same leaf); B separates them by
driver profile, C then names them. So B$\to$C yields
\emph{mechanism-differentiated} zones; C-alone yields intensity bands.

\noindent\textbf{When to use which.} Build C-alone first (cheap). Only reach for
B$\to$C if intensity-only zones visibly lump mechanistically-different regions.
If the modulation is $\sim$1-D (PCA PC1\,$=67\%$ suggests it may be), C-alone
suffices and B adds little. \textbf{Recommendation:} lead with \textbf{C} as the
flagship (portable rule-based zonation); use \textbf{B} to explain what each
zone means; keep \textbf{A} as a throwaway check. Final manuscript object: a
zonation map + a table giving each zone its (i)~defining rule,
(ii)~characteristic $\tau$-modulation, (iii)~dominant drivers.

% =============================================================================
\section{Caveats that must be resolved before drawing final figures}
% =============================================================================

\begin{enumerate}[nosep]
  \item \textbf{OOB understates the signal.} The diagnostics above use OOB, where
  climate $R^2=0.53$; the manuscript uses \emph{spatial} CV, where climate
  $R^2=0.28$. Spatial CV leaves a much larger residual for non-climate to explain,
  so the honest (spatial-CV) effect sizes are \emph{larger} than the
  3.7\%/13.4\% quoted here. \textbf{Re-validate the M5/M7 modulation fields with
  spatial-block CV} (reuse the existing 2$^\circ$ blocks).
  \item \textbf{Edaphic $\Delta R^2$ discrepancy.} Quick OOB gives
  edaphic$+$landuse only $+0.017$, whereas the manuscript reports edaphic-alone
  $+0.076$ (spatial CV). Almost certainly the OOB-vs-spatial-CV difference, but it
  is central to the variance-partition claims and must be reconciled explicitly.
  \item \textbf{Map coherence.} At 0.1$^\circ$ the modulation field is spatially
  noisy, so quantile-bin zones fragment. Coherent (Köppen-like) zones require
  coarser resolution or spatially-constrained clustering.
  \item \textbf{The 29\% may be escapable.} It is the \emph{intersection} of
  $\sim$9 biological layers. A ``biology-light'' model carrying only the single
  strongest biological variable --- fungal proportion (Yu 2022, $\sim$1\,km,
  likely near-global) --- would intersect far less and could push the biological
  tier toward $\sim$49\% while keeping the signal that mattered most in the ALE.
  \textbf{Check per-layer coverage first} (SPUN richness/endemism is the likely
  bottleneck).
  \item \textbf{Biological footprint biased \emph{within} the 29\%.} Microbial
  data are dense in temperate Europe/N.\ America, thin in the tropics and high
  latitudes; biological conclusions are therefore doubly scoped.
\end{enumerate}

% =============================================================================
\section{Where we stand, and the next step}
% =============================================================================

\textbf{Standing:} (i)~storyline~A is chosen and viable; (ii)~the prediction
pipeline is verified sound; (iii)~the global zonality map is achievable as a
continuous modulation field $+$ discrete companion, but only \emph{supervised}
constructions work; (iv)~the abiotic field (M5) is global-but-modest while the
strong signal is biological-but-coverage-limited --- now framed as a finding via
a dual-layer (M5 global $+$ M7 enrichment) deliverable.

\textbf{Immediate next step (on reopen):} set up \textbf{spatial-CV validation}
of the M5 and M7 modulation fields, which simultaneously (a)~yields
paper-consistent effect sizes for the zonality signal and (b)~reconciles the
edaphic $\Delta R^2$ against the manuscript's $+0.076$. The magnitude this returns
decides whether the zonality story is ``modest but real'' or ``substantial,''
i.e.\ how hard we can lean on framing~A.

\textbf{Then:} (a)~check per-layer biological coverage (is the 29\% escapable via
a fungal-proportion-only model?); (b)~build the clustered zone map --- lead with
Route~C (rule-based zonation) at a coherent resolution, add Route~B (TreeSHAP
contributions) to characterise each zone, with Route~A as a sanity check;
(c)~decide whether to pursue microbial gap-filling (Option~1) as a sensitivity
analysis to extend the biological tier globally. The deliverable is a global
Level-1 abiotic zonality map + a Level-2 biological-enrichment zoom-in, each zone
described by rule, $\tau$-modulation, and dominant drivers.

% =============================================================================
\section{Session update (2026-06-26): validation, scale, provenance,
and a data-source correction}
% =============================================================================

\paragraph{Spatial-block CV validation (common row set $n=101{,}400$, 2\textdegree{}
blocks, seed~42, 300 trees).} All quantities computed on a single biological-complete
row set so nested increments are strictly comparable (unlike the manuscript table,
where each model used a different subset). Total $R^2_\text{CV}=0.378$.
\begin{itemize}\itemsep2pt
  \item \textbf{Order-independent attribution (Shapley, sums exactly to total):}
  Climate \textbf{32.8\%}, Edaphic \textbf{32.6\%}, Biological \textbf{24.2\%},
  LandUse 10.4\%. \emph{Climate is only $\sim$1/3 of the attributable signal; the
  non-climate domains carry 2/3.} This is the headline and it is scale-independent.
  \item \textbf{Singles alone:} Edaphic 0.318 $\ge$ Climate 0.306 $>$ Biological
  0.268 $\gg$ LandUse 0.129. Singles sum to 1.021 vs full 0.378 $\Rightarrow$
  overlap 0.64 (``correlated, not redundant''). Use \emph{singles-first} as the
  Results opening.
  \item \textbf{Order-dependence reconciled:} biology added last (CELB$-$CEL)
  $=+0.008$ but Shapley share 24\% --- same data; LandUse redundant with Edaphic.
  \item \textbf{Modulation field is out-of-sample valid but modest:}
  cor(observed climate-residual, full modulation)~$=0.32$ ($R^2=0.10$); discrete
  5--8 zones capture $\eta^2\approx0.07$--$0.09$. Biology is a spatially
  \emph{separate} axis: cor(abiotic modulation, biological enrichment) $=-0.16$.
\end{itemize}

\paragraph{Scale spectrum and noise floor (appendix material).} Residual variogram
nugget/sill $\approx\mathbf{0.66}$: $\sim$2/3 of the beyond-climate residual is
unstructured noise (measurement $+$ change-of-support $+$ sub-reliable-resolution
covariate error) at $>$15~km; the variogram is nearly pure-nugget beyond 15~km.
$\eta^2$/null shows coarse-scale organisation rising from $5\times$ (1\textdegree)
to $33\times$ (20\textdegree) chance, but the absolute coarse signal is small
($\sim$3--7\% of variance), visible only after block-averaging. Covariates predict
the \emph{block-mean} (zonal) residual 2--3$\times$ better than the point residual
($R^2$ 0.10~$\to$~0.20--0.29, peaking at $\sim$3--4\textdegree). \emph{Implication:}
build the map at the coarse effective scale; never show native-pixel texture; the
scale story is supporting (appendix), not the headline. Figure:
\texttt{plots/appendix/scale\_noisefloor.png}.

\paragraph{Provenance control --- PASSES.} Tested whether the coarse signal is a
source/lab artifact using \texttt{source\_db} (44 contributing databases). Only
\textbf{2.4\%} of beyond-climate residual variance is between sources; de-meaning by
source leaves the coarse organisation almost intact ($10$\textdegree: $24\times\to19\times$
chance) and covariate skill barely dented ($0.22\to0.19$ at 5\textdegree). The
variogram mid-lag bump is unchanged by de-meaning, so it is \emph{not} a between-source
offset. Conclusion: the modest zonal signal is \textbf{ecological, not provenance}.
(Caveat: source is confounded with geography, so this is a conservative/upper-bound
test --- surviving signal is definitely ecological.)

\paragraph{Data-source correction (important).} The SOC data is \textbf{not} WoSIS.
It is the OpenLandMap \emph{Open Compendium of Soil Datasets: Soil Observations and
Measurements} (Zenodo~15593990; \texttt{sol\_chem.pnts\_horizons.rds}). WoSIS is not
even a \texttt{source\_db} label in the data. Manuscript fixed in 4 places
(\S Soil data compilation, OC$=$0 sentence, Results data paragraph, Data
availability) $+$ new bib entry \texttt{OpenLandMapSoilDB2025} (author/year
\textbf{TODO}: verify against Zenodo record). Source breakdown corrected: post-QC
198{,}382 profiles from \textbf{44} (not 47) databases; \texttt{GLanCE.points}
(claimed 21.1\% largest) is removed entirely by QC --- the old paragraph applied the
pre-QC breakdown to the post-QC count. Real top sources: USDA\_NCSS 20.4\%,
Netherlands.BIS 11.5\%, SISLAC 8.9\%, CSIRO\_NatSoil 7.0\%, AfSPDB 6.7\%.

\paragraph{Refined storyline.} Lead with \emph{attribution} (climate only 1/3,
robust and scale-free) and the ESM contrast: ESMs get zonality wrong in two
independent ways --- they \emph{over-amplify} the climate range (polar 50--100$\times$
vs data $\sim$3--5$\times$; magnitude claim, caveated by steady-state/change-of-support)
\emph{and} omit the soil$+$biological axis that carries most of the real organisation
(structural claim, robust). Bridge to next-generation models needing explicit
mineral$+$microbial processes. Zonality is \emph{real but modest}; present the map
honestly at coarse scale.

% =============================================================================
\section{Open issues / next steps}
% =============================================================================
\begin{enumerate}\itemsep2pt
  \item \textbf{Promote to production} (must-do before anything enters the MS):
  house-style \texttt{13c\_commonality\_validation.R} at 500 trees, writing
  \texttt{outputs/13c\_*.rds}, the singles$\to$pairs$\to$full $+$ Shapley figure,
  the scale/variogram appendix figure, and the provenance-control table. Repeatable
  via \texttt{run\_pipeline.R}.
  \item \textbf{Biological-parsimony appendix (the ``29\% killer'').} Check
  collinearity + per-layer coverage of the 9 biological layers; single-layer
  marginal $\Delta R^2$ vs coverage; recommend one parsimonious layer to extend the
  biological footprint well beyond 29\%. Figures $\to$ \texttt{plots/appendix/};
  TODO comment already placed in \texttt{manuscript.tex} before
  \texttt{Supplementary tables}.
  \item \textbf{Verify the compendium citation} (authors/year/version) from the
  Zenodo record; update \texttt{OpenLandMapSoilDB2025}.
  \item \textbf{Add Moran's I} on block-aggregated residuals as the standard
  spatial-autocorrelation statistic to accompany the $\eta^2$/null backbone.
  \item \textbf{Build the dual-layer zonality map} at the coarse effective scale
  ($\sim$3--4\textdegree), with per-zone driver characterisation (ALE/TreeSHAP).
  \item \textbf{Reconcile MS magnitude numbers}: the current ``$\sim$10~pp beyond
  climate'' is partly a subset artifact; on comparable rows climate~$=0.31$ and
  beyond-climate~$\approx$7~pp.
  \item \textbf{Narrative}: write with Schimel OCAR; keep Jungian archetypes as a
  hidden internal scaffold only (deferred --- do not implement in text).
\end{enumerate}

\section{Session update (2026-06-26, cont.): production decomposition,
biological parsimony, Tier-1 Moran's I}

\paragraph{Validation promoted to house-style code.} The exploratory spatial-CV
scratchpad is now the production script
\texttt{13c\_commonality\_validation.R} (500 trees, 2\textdegree{}/10-fold/seed-42
block CV; env knobs \texttt{MRT\_13C\_BIO}, \texttt{MRT\_13C\_BLOCK}). It fits all
15 group coalitions on \emph{one common-row set} so coalition $R^2$ are comparable
and Shapley values sum exactly to the full-model $R^2$. The all-9 sensitivity run
reproduces the 300-tree memo numbers almost exactly (Climate 32.9, Edaphic 32.4,
Biological 24.3, LandUse 10.4\,\%; full $R^2=0.378$; $n=101{,}400$), validating the
machinery.

\paragraph{Biological predictor decision (drop Barcel\'o root colonization).} Per-layer
\emph{observation} coverage is not uniform: \texttt{fungal\_proportion} 99.8\,\%,
the five SPUN richness/endemism layers 96.2\,\%, but the three Barcel\'o
root-colonization layers only $\sim$62\,\% --- the sole constraint collapsing the
common-row set to $n=101{,}400$. Equal-row block-CV (script
\texttt{13d\_biological\_parsimony.R}) gives a unique Barcel\'o-3 contribution of
$+0.001$ $R^2$ (inert; alone $-0.004$). \textbf{MAIN model} therefore uses
\texttt{fungal\_proportion} + 5 SPUN layers (6 layers, 96\,\% coverage, $\sim$84\,\%
of the biological $\Delta R^2$); the all-9 set is retained as a sensitivity case.
Collinearity: only the derived ratios are redundant
(\texttt{EcM\_AM\_root\_ratio}$\approx$\texttt{EcM\_roots} $0.99$;
\texttt{EcM\_AM\_richness\_ratio}$\approx$\texttt{EcM\_richness} $0.92$); biology is
otherwise multidimensional, so no single layer substitutes for the set.

\paragraph{Headline decomposition (MAIN, $n=154{,}675$, full $R^2=0.380$).} Dropping
Barcel\'o lifts the common-row set by 53\,\% with essentially no loss of full-model
skill (a robustness point). Shapley shares: \textbf{Climate 35.1, Edaphic 32.8,
Biological 22.3, LandUse 9.7\,\%} (sum exact) --- climate $\approx\!\tfrac13$,
non-climate $\approx\!\tfrac23$. Singles: $C\,0.307 \approx E\,0.307 > B\,0.227
\gg L\,0.110$; $\Sigma$ singles $=0.950$ vs full $0.380 \Rightarrow$ overlap $0.57$
(``correlated, not redundant''). Land-use unique $\approx 0$; biology is badly
\emph{understated} by its nested-last increment ($+0.010$) versus Shapley
($22\,\%$) --- \textbf{lead Results with Shapley}. Conclusion is insensitive to the
biological-set choice (vs all-9: 32.9/32.4/24.3/10.4). Figure:
\texttt{plots/step\_13c\_commonality/13c\_decomposition.png}.

\paragraph{Tier-1 rigour --- Moran's I (script \texttt{13e\_tier1\_morans\_i.R}).} The
beyond-climate residual (observed $-$ climate-only block-CV prediction), aggregated
to spatial blocks, has \textbf{Moran's $I=0.25$ at 2\textdegree{} and $0.26$ at
5\textdegree{}} ($p \ll 0.001$): the residual is spatially \emph{organised}, not
noise, complementing the weak/scale-dependent ICC. (Cross-check: the climate-only
block-CV $R^2=0.307$ here matches the 13c climate single exactly.) Moderate
magnitude is consistent with the noise floor; zonality is real but modest.

\paragraph{Manuscript reconciliation (analysis done; prose pending discussion).} The
current ``$\sim$10 percentage points beyond climate'' (L826--827) is a
\emph{footprint artifact}: climate $R^2=0.280$ is scored on $n\approx178$k while the
full model is on $n\approx101$k. On a common footprint climate rises to $0.307$ and
the increment is \textbf{$\sim$7 pp}. Full reconciliation table (old$\to$new with
line refs; mechanical fixes vs.\ framing decisions) in
\texttt{results\_reconciliation\_2026-06-26.md}. The qualitative ordering (edaphic
largest non-climate, biology second, land-use redundant) survives.

\paragraph{Next.} Run the 1\textdegree{}/5\textdegree{} block-size sweep
(\texttt{MRT\_13C\_BLOCK}) for Shapley stability; then the Results-narrative
discussion (Shapley-primary framing; replace/augment the M1--M7 table) and the
manuscript edits.

\bigskip
\noindent\rule{\linewidth}{0.4pt}\\
{\footnotesize Artifacts (2026-06-25/26):
\texttt{outputs/MRT\_predictions/zonality\_modulation\_M5\_M1\_logratio.tif};
\texttt{plots/appendix/scale\_noisefloor.png}; exploratory scripts in the session
scratchpad (\texttt{spatialcv\_validation.R}, \texttt{spatialcv\_finish.R},
\texttt{scale\_noisefloor.R}, \texttt{provenance\_control.R}; earlier
\texttt{zone\_explore.R}, \texttt{zone\_supervised.R}, \texttt{m5\_zonality.R}).
All reproducible from \texttt{outputs/12b\_model\_ready.rds},
\texttt{outputs/13\_rf\_models.rds}, \texttt{outputs/13\_var\_groups.rds}.
\emph{Scratchpad scripts are exploratory and must be promoted to house-style
pipeline code before any result enters the manuscript.}}

\end{document}
